% !TeX root = ../../Book2.tex
% 纲要标题：### 17.1 单代数与中心
% 纲要位置：计划纲要.md，第 1596 行。

\section{单代数与中心}
\label{b2:sec:1701}

半单结构定理把有限维代数中的单块写成除环上的矩阵环，但底域未必就是这些块的中心。中心中的元素同全部乘法相容，所以每个单块也自然成为其中心上的代数。研究一个单块扩大标量后的结构时，就要先确定这个标量域：中心单代数要求它恰为全部中心，接下来将证明这一条件使任意扩域后的代数仍然单。张量积和中心化子的结构也将以此为基础。本章的代数结合且含单位元，子代数与嵌入保持单位元；中心单代数一律非零且有限维。

% 纲要要点（供目录核对）：
%
% * 有限维中心单代数；固定底域与实际中心，包络判据核对稠密性所用右除环及全线性算子目标
% * 分裂代数与分裂域；区别给定映射的同构检测与扩域后新同构的存在，核对有限维自同态的标量扩张、基比较下降及次数平方，正式证明纯不可分扩域后的非单性反例
% * 张量积与反代数；用两个包络作用证明张量积中心单，区别一般单代数的非单张量积
%

\subsection{有限维中心单代数}
\label{b2:subsec:1701:01}

\begin{definition}\label{b2:def:central-simple-algebra}
设 \(k\) 为域，\(A\) 为非零有限维含幺结合 \(k\) 代数。若 \(A\) 没有非零真双边理想，且其中心满足 \(Z(A)=k1_A\)，则称 \(A\) 为 \(k\) 上的\term[zhongxin-dan-daishu]{中心单代数}[central simple algebra]。下文用结构嵌入把 \(k1_A\) 识别为 \(k\)。
\end{definition}

单性限制的是双边理想，中心性还要求指定的底域就是全部中心。例如，有限域扩张 \(L/k\) 作为 \(k\) 代数是单的，但中心为 \(L\)，只有 \(L=k\) 时才中心单。对 \(n\ge1\)，矩阵代数 \(M_n(k)\) 则中心单：单性由\cref{b2:prop:simple-artinian}证明，中心为标量矩阵是\cref{b2:prop:matrix-centers-central-idempotents}在除环为 \(k\) 时的结论。

\ProseParagraphBreak
底域的选择确实会改变这个判断。例如 \(M_2(\mathbb C)\) 作为实代数是单的，中心却是 \(\mathbb C I_2\)，大于 \(\mathbb R I_2\)；它的实维数为 \(2\cdot2^2=8\)。改以 \(\mathbb C\) 为底域后，同一个乘法代数成为中心单代数，维数为 \(2^2=4\)。因此“中心单”要连同底域一起说明；后面得到的维数平方公式也只对这个指定的中心域成立。

\ProseParagraphBreak
一般非零含幺单代数 \(B\) 的中心是域，这不需要有限维假设。若 \(0\ne z\in Z(B)\)，则 \(Bz\) 是非零双边理想，故等于 \(B\)，从而 \(bz=1\) 对某个 \(b\) 成立。由 \(z\) 中心可知 \(zb=1\)，且从 \(zx=xz\) 可推出 \(z^{-1}x=xz^{-1}\)。所以 \(z\) 在中心内可逆。若 \(B\) 还是有限维 \(k\) 代数，\(Z(B)\) 作为 \(B\) 的 \(k\) 子空间也有限维，因而是 \(k\) 的有限扩张。

\ProseParagraphBreak
Artin–Wedderburn 定理给出 \(A\cong M_r(D)\) 时，除环 \(D\) 的中心未必就是原底域 \(k\)。中心单条件进一步固定这个参数，要求 \(Z(D)=k\)。

\begin{proposition}\label{b2:prop:csa-division-description}
设 \(k\) 为域，\(A\) 为中心单 \(k\) 代数，则存在整数 \(r\geq1\) 与有限维中心除 \(k\) 代数 \(D\)，使 \(A\cong M_r(D)\)。反过来，对任意这样的 \(r,D\)，矩阵代数 \(M_r(D)\) 中心单。整数 \(r\) 与 \(D\) 的 \(k\) 代数同构类型由 \(A\) 唯一确定。
\end{proposition}
\begin{proof}
有限维使左理想降链稳定，单性与\cref{b2:prop:simple-artinian}给出 \(A\cong M_r(D)\)。这个构造保持中心中的 \(k\) 标量，因此是 \(k\) 代数同构，\(D\) 有限维。由\cref{b2:prop:matrix-centers-central-idempotents}，\(Z\bigl(M_r(D)\bigr)=Z(D)I_r\)，而 \(A\) 的中心恰为 \(k\)，于是 \(Z(D)=k\)。反向的单性由\cref{b2:prop:simple-artinian}中矩阵环的结论得到，中心由同一中心公式给出。唯一性是\cref{b2:prop:semisimple-multiplicity}对正则模及其自同态除环的唯一性，构造同时保留 \(k\) 作用。
\end{proof}

\begin{theorem}[包络代数判据]\label{b2:thm:csa-enveloping-isomorphism}
设 \(k\) 为域，\(A\ne0\) 为有限维含幺结合 \(k\) 代数，则 \(A\) 中心单，当且仅当左右乘作用给出同构
\begin{equation}\label{b2:eq:csa-enveloping-isomorphism}
\Theta_A:A\otimes_kA^{\mathrm{op}}\longrightarrow\operatorname{End}_k(A),
\qquad a\otimes b^{\mathrm{op}}\longmapsto(x\mapsto axb).
\end{equation}
\end{theorem}
\begin{proof}
\cref{b2:thm:enveloping-bimodules}已证明 \(\Theta_A\) 是代数同态。若 \(A\) 中心单，要证明它满射，就需让左右乘的有限和实现任意 \(k\) 线性算子。把正则双模 \(A\) 看成左 \(A^e\) 模，其子模正是双边理想，故这个模简单。由\cref{b2:prop:regular-bimodule-endomorphism}，稠密性定理所用的右标量除环为
\[
D=\operatorname{End}_{A^e}(A)^{\mathrm{op}}=Z(A)^{\mathrm{op}}=k,
\]
所以目标 \(\operatorname{End}_D(A_D)\) 恰为 \(\operatorname{End}_k(A)\)。\(A\) 对这个右除环有限维，\cref{b2:cor:density-finite-dimensional}用于环 \(A^e\) 及简单模 \(A\)，便使 \(\Theta_A\) 满射。来源与目标的 \(k\) 维数同为 \((\dim_kA)^2\)，故它还是单射。

反过来，若 \(\Theta_A\) 同构，则每个双边理想 \(I\) 都是正则双模 \(A\) 的 \(A^e\) 子模，因而被 \(\Theta_A\) 的全部作用算子保持。满射性使它被每个 \(T\in\operatorname{End}_k(A)\) 保持。若 \(0\ne x\in I\)，对任意 \(y\in A\)，把 \(x\) 扩充为一组基并指定 \(T(x)=y\)，就得到 \(y=T(x)\in I\)。所以 \(I=A\)，证明单性。由正则双模自同态公式，\(Z(A)\) 同构于 \(\operatorname{End}_{A^e}(A)\)；后者现在是全部线性算子的中心化子。在\cref{b2:prop:tensor-factor-centralizers}中取 \(U=A\)、\(W=k\)，便知这个中心化子为 \(k\operatorname{id}_A\)，故 \(Z(A)=k\)。
\end{proof}

张量积中的元素是有限和，所以这个同构具体断言：对每个 \(T\in\operatorname{End}_k(A)\)，都能找到有限组 \(a_\nu,b_\nu\in A\)，使
\[
T(x)=\sum_{\nu=1}^{r}a_\nu x b_\nu\qquad(x\in A).
\]
纯张量 \(a\otimes b^{\mathrm{op}}\) 对应一次左右乘；一般线性算子由这样的操作的有限和实现。

\ProseParagraphBreak
在 \(A=M_n(k)\) 中，若 \(X=(x_{pq})\)，矩阵单位给出
\[
E_{ij}XE_{\ell m}=x_{j\ell}E_{im}.
\]
这项操作读取第 \((j,\ell)\) 个条目，再把它放到第 \((i,m)\) 个位置，其余条目清零。让两组位置遍历所有可能，就得到矩阵空间的全部线性自同态的矩阵单位；它们的线性组合正是任意线性操作。扩大底域后，这些读取、放置公式依然成立，这也显示了包络代数判据适合处理标量扩张的原因。

\subsection{分裂代数与分裂域}
\label{b2:subsec:1701:02}

对域扩张 \(K/k\)，记 \(A_K=K\otimes_kA\)。其乘法逐因子定义，单位为 \(1\otimes1\)。
\begin{definition}\label{b2:def:csa-splitting-field}
设 \(k\) 为域，\(A\) 为中心单 \(k\) 代数。若存在整数 \(n\geq1\) 使 \(A\cong M_n(k)\)，则称 \(A\) 为\term[fenlie-daishu]{分裂代数}[split algebra]，或称它在 \(k\) 上分裂。对域扩张 \(K/k\)，记 \(A_K=K\otimes_kA\)；若 \(A_K\cong M_n(K)\)，则称 \(K\) 为 \(A\) 的\term[fenlie-yu-daishu]{分裂域}[splitting field of an algebra]。
\end{definition}
\symidx{\(A_K=K\otimes_kA\)：代数的标量扩张}

扩张标量能检测一个给定映射是否为同构，但这与扩张后才找到一个同构不同。若 \(f:V\to W\) 是给定的 \(k\) 线性映射，且 \(K\otimes_k f\) 为同构，则 \(f\) 本身就是同构：向量空间的短正合列分裂，张量后仍正合，所以 \(f\) 的核与余核扩张后都为零；非零 \(k\) 向量空间的一组基扩张后仍非空，因而核与余核原本就为零。对于给定的 \(k\) 代数同态，这同样检测它是否为代数同构。

\ProseParagraphBreak
若只知道 \(A_K\cong B_K\)，这个同构却未必来自某个 \(k\) 代数同态 \(A\to B\)。例如实四元数代数 \(\mathbb H\) 与 \(M_2(\mathbb R)\) 扩张到 \(\mathbb C\) 后都同构于 \(M_2(\mathbb C)\)，其中四元数的分裂将在\cref{b2:prop:quaternion-basis-splitting}中具体给出；但 \(\mathbb H\) 是除代数，\(M_2(\mathbb R)\) 有非零不可逆矩阵，所以两者在 \(\mathbb R\) 上不同构。\cref{b2:cor:similarity-descent}中矩阵相似能够下降，是因为首一不变因子在扩张标量后保留，这个额外的分类信息不能直接推广到任意代数的同构问题。

\begin{theorem}[标量扩张]\label{b2:thm:csa-base-change}
设 \(k\) 为域，\(A\) 为中心单 \(k\) 代数，则对任意域扩张 \(K/k\)，标量扩张 \(A_K=K\otimes_kA\) 为中心单 \(K\) 代数，不要求扩张有限或可分。
\end{theorem}
\begin{proof}
要保留中心单性，可以把包络代数判据中的同构扩张标量。由于 \(A\) 有限维，取它的一组有限基后，矩阵单位给出自然同构
\[
K\otimes_k\operatorname{End}_k(A)\cong\operatorname{End}_K(A_K),
\]
将 \(c\otimes T\) 送到 \(d\otimes a\mapsto cd\otimes T(a)\)。它把两侧相应的有限矩阵单位基一一对应，所以是代数同构；这个与自同态代数交换的公式不能无条件用于无限维空间。另有
\[
A_K\otimes_K(A_K)^{\mathrm{op}}
\cong K\otimes_k(A\otimes_kA^{\mathrm{op}}).
\]
将\eqref{b2:eq:csa-enveloping-isomorphism}扩张标量并作这些识别，所得映射正是 \(A_K\) 的左右乘作用。它是同构，且 \(A_K\ne0\)，因为原 \(k\) 基扩张后仍为 \(K\) 基。因此包络代数判据给出所需中心单性。
\end{proof}

若只假设 \(A\) 单，标量扩张后却可能出现非零真双边理想。\cref{b2:prop:inseparable-base-change-counterexample}将给出一个有限纯不可分域扩张作为 \(k\) 代数的反例；它的中心大于 \(k\)，因而不满足上述定理的中心性条件。

\begin{proposition}[中心单性的下降]\label{b2:prop:csa-descent-from-extension}
设 \(k\) 为域，\(A\ne0\) 为有限维含幺结合 \(k\) 代数，\(K/k\) 为域扩张。若 \(A_K=K\otimes_kA\) 为中心单 \(K\) 代数，则 \(A\) 为中心单 \(k\) 代数。
\end{proposition}
\begin{proof}
若 \(I\) 是 \(A\) 的非零真双边理想，则 \(K\otimes_kI\) 是 \(A_K\) 的双边理想。取 \(I\) 的一组 \(k\) 基并扩充为 \(A\) 的基，便知扩张后它仍非零且仍为真子空间，这与 \(A_K\) 单矛盾。因此 \(A\) 单。

再把 \(1_A\) 扩充为 \(A\) 的 \(k\) 基。若 \(z\in Z(A)\)，则 \(1\otimes z\) 与 \(K\) 标量及全部 \(1\otimes a\) 交换，所以在 \(Z(A_K)=K(1\otimes1_A)\) 中。于是 \(1\otimes z=c\otimes1_A\) 对某个 \(c\in K\) 成立。用扩张后的基比较系数，\(z\) 的所有非单位基坐标都为零，其单位坐标原本在 \(k\) 中，故 \(z\in k1_A\)。这证明中心性。
\end{proof}

\begin{corollary}\label{b2:cor:csa-algebraic-closure-splits}
设 \(k\) 为域，\(\Omega\) 为其代数闭包，\(A\) 为中心单 \(k\) 代数，则 \(A_\Omega=\Omega\otimes_kA\cong M_n(\Omega)\)，其中整数 \(n\geq1\) 唯一，且 \(\dim_kA=n^2\)。
\end{corollary}
\begin{proof}
\cref{b2:thm:csa-base-change}使 \(A_\Omega\) 有限维且中心单，因而简单、半单。\cref{b2:prop:finite-dimensional-semisimple-algebra}在代数闭域上把它写成矩阵环的有限乘积；单性排除多个非零因子，所以只剩 \(M_n(\Omega)\)。扩张标量不改变有限基的大小，故 \(\dim_kA=n^2\)，这也唯一确定 \(n\)。
\end{proof}

称 \(n=\sqrt{\dim_kA}\) 为 \(A\) 的\term[zhongxin-dan-daishu-cishu]{次数}[degree of a central simple algebra]，记为 \(\deg A\)。次数为 \(n\)，整个代数的向量空间维数则为 \(n^2\)；因此中心单代数的维数必须是完全平方，例如三维或六维的代数不能中心单。任意分裂域给出的矩阵大小都为这个 \(n\)，因为维数在标量扩张下不变。\secref{b2:sec:1702}还将证明可以选取有限 Galois 分裂域；仅知道在代数闭包上分裂，尚未说明有限可分扩张的存在。
\symidx{\(\deg A=\sqrt{\dim_kA}\)：中心单代数的次数}

\subsection{张量积与反代数}
\label{b2:subsec:1701:03}

反代数 \(A^{\mathrm{op}}\) 与 \(A\) 有相同的中心，双边理想也相同，所以中心单性在取反代数后保持。它的次数等于 \(\deg A\)，因为维数不变。与自身的反代数作张量则有更强的结论：由\cref{b2:thm:csa-enveloping-isomorphism}，
\[
A\otimes_kA^{\mathrm{op}}\cong M_{\dim_kA}(k)
=M_{(\deg A)^2}(k).
\]
无论 \(A\) 本身是否分裂，这个张量积总是分裂的。右侧矩阵的阶数是 \(\dim_kA=(\deg A)^2\)，不能误记为 \(\deg A\)。在\secref{b2:sec:1801}的 Brauer 群中，分裂矩阵代数代表单位元，这个同构将使 \(A^{\mathrm{op}}\) 的类成为 \(A\) 的类的逆元。

\begin{theorem}[中心单代数的张量积]\label{b2:thm:csa-tensor-product}
设 \(k\) 为域，\(A,B\) 为中心单 \(k\) 代数，则 \(A\otimes_kB\) 也是中心单 \(k\) 代数，且
\(\deg(A\otimes_kB)=\deg A\deg B\)。
\end{theorem}
\begin{proof}
令 \(C=A\otimes_kB\)。\(C\) 的左右乘可以分别在 \(A,B\) 两个因子上实现，因此要把它的包络作用与两个因子的包络作用相比较。交换不同张量因子的次序并取结合识别，得到
\[
C\otimes_kC^{\mathrm{op}}
\cong(A\otimes_kA^{\mathrm{op}})\otimes_k(B\otimes_kB^{\mathrm{op}}).
\]
分别应用两个包络代数同构，再用\cref{b2:prop:tensor-factor-centralizers}中已经通过矩阵单位证明的同构，得到
\[
C\otimes_kC^{\mathrm{op}}
\cong\operatorname{End}_k(A)\otimes_k\operatorname{End}_k(B)
\cong\operatorname{End}_k(C).
\]
在纯张量上，这个复合把 \((a\otimes b)\otimes(a'\otimes b')^{\mathrm{op}}\) 送到
\(x\otimes y\mapsto axa'\otimes byb'\)，正是 \(C\) 的左右乘作用。因此包络代数判据保证 \(C\) 中心单。维数乘法给出次数乘法。
\end{proof}

特别地，\(M_r(A)\cong M_r(k)\otimes_kA\) 是中心单代数，次数为 \(r\deg A\)。同构把 \(E_{ij}\otimes a\) 送到第 \((i,j)\) 个条目为 \(a\)、其余为零的矩阵；基与乘法同时对应。另一方面，两个一般单 \(k\) 代数的张量积不必单，中心性在定理中不能删除。例如
\[
\mathbb C\otimes_{\mathbb R}\mathbb C\longrightarrow\mathbb C\times\mathbb C,
\qquad z\otimes w\longmapsto(zw,z\bar w)
\]
是实代数同构。确实，按第一个 \(\mathbb C\) 因子看，来源以 \(1\otimes1,1\otimes i\) 为基，两个像为 \((1,1),(i,-i)\)，线性无关且生成目标；乘法直接保持。恒等嵌入和复共轭使第二因子分别产生两个坐标，而目标中的 \(\mathbb C\times0\) 是非零真双边理想，故张量积不是单环。这里每个因子作为实代数都单，但其中心为 \(\mathbb C\)，大于底域 \(\mathbb R\)。

\ProseParagraphBreak
把底域换成交换环时，有限维性应改为忠实有限投射性；左右乘作用的同构则仍有意义。\appref{b2:app:C}的\cref{b2:def:azumaya-algebra}据此定义 Azumaya 代数\footnote{東屋五郎（假名：あずまや ごろう，罗马音：Gorō Azumaya，1920-2010），日本数学家，研究环论及代数的结构；Azumaya 代数以他命名。}，并由\cref{b2:thm:azumaya-fiber-criterion}证明其剩余域纤维判据。这项推广不将每个纤维预设为已经分裂的矩阵代数。

\begin{proposition}[不可分标量扩张后的非单性]\label{b2:prop:inseparable-base-change-counterexample}
设 \(p\) 为素数、\(t\) 为 \(\mathbb F_p\) 上的超越元，令 \(k=\mathbb F_p(t)\)。在 \(k\) 的代数闭包中取 \(\alpha\) 满足 \(\alpha^p=t\)，并令 \(L=k(\alpha)\)。则 \(X^p-t\) 在 \(k[X]\) 中不可约，\([L:k]=p\)，且以第一个张量因子为 \(L\) 标量时有含幺 \(L\) 代数同构
\[
L\otimes_kL\cong L[\varepsilon]/(\varepsilon^p).
\]
其中 \(1\otimes\alpha-\alpha\otimes1\) 对应 \(\varepsilon\) 的剩余类。它是非零幂零元，生成一个非零真理想，故 \(L\otimes_kL\) 不是单代数。原代数 \(L\) 虽为单 \(k\) 代数，但其中心为 \(L\ne k\)。
\end{proposition}
\begin{proof}
先证明 \(t\) 在 \(k\) 中不是 \(p\) 次方。若 \((f/g)^p=t\)，其中 \(f,g\in\mathbb F_p[t]\)、\(g\ne0\)，约去公因子后有 \(f^p=tg^p\)。左边不可约因子 \(t\) 的指数为 \(p\) 的倍数，右边的指数则模 \(p\) 余一，矛盾。

令 \(\mu_{\alpha,k}\) 为 \(\alpha\) 在 \(k\) 上的首一最小多项式。它整除 \(X^p-t\)，而在代数闭包中 \(X^p-t=(X-\alpha)^p\)，故 \(\mu_{\alpha,k}=(X-\alpha)^d\)，其中 \(1\le d\le p\)。若 \(d<p\)，其 \(X^{d-1}\) 系数为 \(-d\alpha\in k\)。此时 \(d\) 在特征 \(p\) 的域中非零，便有 \(\alpha\in k\)，与 \(t\) 不是 \(p\) 次方矛盾。因此 \(d=p\)，\(\mu_{\alpha,k}=X^p-t\)，并且 \(1,\alpha,\ldots,\alpha^{p-1}\) 是 \(L/k\) 的基。

把 \(L[X]\) 中的 \(X\) 送到 \(1\otimes\alpha\)，把系数 \(c\in L\) 送到 \(c\otimes1\)，得到含幺 \(L\) 代数同态。关系 \(X^p-t\) 在像中为零，故它诱导
\[
L[X]/(X^p-t)\longrightarrow L\otimes_kL.
\]
两侧分别以 \(X^j\) 的剩余类和 \(1\otimes\alpha^j\)（\(0\le j<p\)）为 \(L\) 基，且这些基向量一一对应，所以此映射为同构。在 \(L[X]\) 中令 \(\varepsilon=X-\alpha\)，便有 \(X^p-t=\varepsilon^p\)，得到所述截断多项式环。它以 \(1,\varepsilon,\ldots,\varepsilon^{p-1}\) 的剩余类为 \(L\) 基，故 \(\varepsilon\ne0\) 而 \(\varepsilon^p=0\)；理想 \((\varepsilon)\) 非零，且商为非零域 \(L\)，所以是真理想。最后，\(L\) 是域，故作为 \(k\) 代数单；但 \([L:k]=p>1\)，且 \(Z(L)=L\)，因而它不中心单，也就不满足\cref{b2:thm:csa-base-change}的假设。
\end{proof}
